\section{The pure fragment}\label{sec:pure}

For the exact pop counts and ordering formulas in this section, restrict attention to closed numerical terms over deterministic, total computable operators. Thus every tuple of numerical operands is admissible for its operator. In the broader fragment of the introduction, operators may be partial and an inert circle may remain: every reduction still has at most $\size{t}$ pops, but reaching a numeral in exactly that many pops requires admissibility throughout. A ready circle is eligible for a pop only when its operand tuple is admissible.

\begin{definition}
A \emph{pure term} is defined inductively: a numeral is a pure term, and if $o$ is an operator of arity $k$ and $t_1,\dots,t_k$ are pure terms then the circle $o(t_1,\dots,t_k)$ is a pure term. A circle is \emph{ready} when every operand is a numeral. A \emph{pop} replaces a ready circle by the numeral that its operator computes from its operands. The \emph{size} $\size{t}$ is the number of circles in $t$.
\end{definition}

\begin{proposition}[Termination in exactly $\size{t}$ pops]\label{prop:pure-term}
Every pop strictly decreases $\size{t}$ by exactly one. Consequently every maximal sequence of pops from a pure term $t$ ends, after exactly $\size{t}$ pops, in a numeral, and the numeral does not depend on the order of pops.
\end{proposition}

\begin{proof}
A pop removes one circle, the ready one, and leaves every other circle in place, so the size drops by one. A term with no circles is a numeral, and a term with a circle has a ready circle (take a circle of greatest depth, all of whose operands are then numerals), so the process continues until size zero. For the independence of the order, suppose two distinct circles $c_1$ and $c_2$ are ready in $t$. They are disjoint, since a ready circle has only numerals inside it, and popping one does not alter the other or its operands, so the two pops commute and yield the same term. A standard argument on the number of pops (the diamond property for a terminating system with commuting local steps, Newman's lemma in its simplest form) then gives that all maximal pop sequences end in the same numeral.
\end{proof}

The proposition establishes termination for this fragment, and the extension of Section~\ref{sec:application} is judged against it. Values and predicates can depend on numerical inputs, but the fixed expression supplies no internal unbounded iteration or search. Pop counts measure invocations of the chosen primitives, not their internal arithmetic cost.

\subsection{Counting evaluation orders}

Although the result of a pure term is fixed, the set of ways to reach it is not. A \emph{pop sequence} is a total order on the circles of $t$ in which every circle occurs after all circles in its subtree. These are the linear extensions of the tree order in which a parent is above its children. For a rooted forest the number of linear extensions has a closed form, due to Knuth \cite{knuth} (see also \cite{stanley}), which is the hook length formula for trees.

\begin{theorem}[Hook length for pop orderings]\label{thm:hook}
Let $t$ be a pure term with $n=\size{t}$ circles and let $h(v)$ denote the number of circles in the subtree rooted at the circle $v$. The number of pop sequences of $t$ is
\[
\frac{n!}{\prod_{v} h(v)}.
\]
\end{theorem}

\begin{proof}
Consider a uniformly random total order of the $n$ circles. For a fixed circle $v$, the probability that $v$ occurs after every other circle of its subtree is $1/h(v)$, because among the $h(v)$ circles of the subtree each is equally likely to come last. These events, over the different circles $v$, are independent: the event for $v$ depends only on the relative order of the circles in the subtree of $v$, and conditioned on the relative order inside each of the subtrees of the children of $v$, the position of $v$ among those circles is uniform and independent of the internal orders. By induction on the tree the joint probability that every circle follows its subtree is $\prod_v 1/h(v)$. Multiplying by $n!$ total orders gives the count of valid ones.
\end{proof}

A second count is the number of distinct \emph{states} reachable by pops, that is, the number of distinct partial results, which is the number of down-closed sets of circles (sets closed under taking subtrees).

\begin{proposition}[States as order ideals]\label{prop:ideals}
Let $I(v)$ be the number of sets of circles in the subtree of $v$ that are closed under taking descendants. Then $I(v)=1+\prod_{c}I(c)$, where $c$ ranges over the circle children of $v$, and the empty product is $1$.
\end{proposition}

\begin{proof}
A closed set either omits $v$, in which case it is an arbitrary choice of a closed set in each child subtree, giving $\prod_c I(c)$ possibilities, or contains $v$, in which case closure forces it to contain the whole subtree of $v$, giving one possibility.
\end{proof}

\begin{example}
For the balanced tree with seven circles (a root, two children, each with two circle children), $h$ takes the values $7,3,3,1,1,1,1$ so Theorem~\ref{thm:hook} gives $7!/63=\factBalancedOrderings$ pop sequences. The recursion of Proposition~\ref{prop:ideals} gives $I=2$ for a bottom circle, $I=1+2\cdot 2=5$ for a middle circle, and $I=1+5\cdot 5=\factBalancedIdeals$ for the root. Both counts were verified by memoized dynamic programming over states, against the closed form, on \factPureTrees{} random trees.
\end{example}

The two counts illustrate a contrast that returns in Section~\ref{sec:dp}. The number of paths through the reduction graph is the larger quantity ($\factBalancedOrderings$ here), and the number of states is the smaller ($\factBalancedIdeals$ here). A dynamic program over states computes a quantity indexed by paths in time proportional to the number of states, and the saving, which is exponential in general, is precisely the saving that sharing provides in graph reduction.
